\chapter{INTRODUCTION}

\section{Preamble}
In the modern landscape of software engineering and computer science, data 
structures play a fundamental role in the efficient organization, processing, 
and retrieval of information. Among the diverse array of available data 
structures, the dynamic array stands out as one of the most ubiquitous and 
critical tools in a programmer's arsenal. Unlike static arrays, which require 
the size of the collection to be known at compile time or at the time of 
allocation, dynamic arrays seamlessly adjust their capacity during runtime to 
accommodate an arbitrary number of elements. This flexibility, coupled with 
the ability to provide $O(1)$ constant-time random access, has cemented the 
dynamic array as the default sequential container in almost all major high-level 
programming languages, such as the \texttt{list} in Python, \texttt{std::vector} 
in C++, and \texttt{ArrayList} in Java.

The physical architecture of modern Random Access Memory (RAM) dictates that 
data should ideally be stored contiguously to exploit hardware-level caching 
mechanisms. When elements are stored in adjacent memory locations, the CPU 
can prefetch data into L1, L2, and L3 caches, dramatically reducing latency 
during sequential iterations. The dynamic array achieves this memory-level 
locality by maintaining its elements in a single contiguous block of physical 
memory. However, this contiguity introduces a profound challenge: because memory 
address space is shared among multiple processes and allocations, an array 
cannot simply expand in-place if the adjacent address space is already occupied 
by other data. Consequently, dynamic array implementations must allocate a new, 
larger contiguous block elsewhere in memory and transfer the existing elements 
to this new location when capacity is exceeded.

The pervasive use of dynamic arrays implies that even marginal improvements in 
their underlying implementation can lead to significant macroscopic gains in the 
performance of software applications. These performance benefits range from 
resource-constrained Internet of Things (IoT) devices, where memory is extremely 
scarce, to large-scale distributed systems and data analytics platforms, where 
millions of reallocations occur every second. The core mechanism that enables 
the flexibility of a dynamic array is its growth strategy: when the array reaches 
its capacity, a larger block of memory is allocated, the existing elements are 
copied to the new block, and the old block is deallocated. The efficiency of 
a dynamic array is thus heavily dictated by how it manages this reallocation 
process. If reallocations occur too frequently, the application incurs significant 
computational overhead due to element copying, thereby degrading throughput. 
Conversely, if the array requests excessively large blocks of memory during 
reallocation to avoid future copying, it risks squandering valuable system 
memory, leading to internal fragmentation and potentially degrading system-wide 
performance due to cache misses and swapping.

Achieving an optimal balance between time efficiency (computational overhead) 
and space efficiency (memory waste) is the central challenge in designing 
dynamic array growth algorithms. As software systems increasingly handle 
unpredictable, bursty, and massive volumes of data, the need for intelligent, 
adaptive, and highly efficient memory allocation mechanisms becomes ever more 
pronounced.

\section{Background}
Dynamic arrays operate by maintaining a contiguous block of memory with a 
certain \textit{capacity}, while keeping track of the actual number of elements 
it currently holds, referred to as its \textit{size} or \textit{length}. As long 
as the size is less than the capacity, inserting a new element at the end of 
the array (often called a \texttt{push} or \texttt{append} operation) is a 
trivial $O(1)$ operation, involving only a memory write and an increment of 
the size counter. The critical operational phase occurs when an insertion is 
requested but the array is full (i.e., $\text{size} = \text{capacity}$). At 
this juncture, the dynamic array must undergo a resizing operation.

The standard procedure for resizing involves three primary steps:
\begin{enumerate}
    \item \textbf{Allocation:} Requesting a new, larger contiguous block of 
    memory from the operating system or the language's memory allocator.
    \item \textbf{Copying:} Transferring all existing elements from the old 
    memory block to the newly allocated block.
    \item \textbf{Deallocation:} Freeing the old memory block so it can be 
    reclaimed by the system.
\end{enumerate}

\begin{figure}[H]
\centering
\begin{tikzpicture}[
    scale=0.88, transform shape,
    cell/.style={draw, minimum size=8mm, font=\ttfamily, fill=blue!10},
    empty/.style={draw, minimum size=8mm, font=\ttfamily, fill=gray!10},
    moved/.style={draw, minimum size=8mm, font=\ttfamily, fill=green!15},
    arrow/.style={-Stealth, thick}
]
    % Old Array
    \node[anchor=west] (o_title) at (-4.8, 1.0) {\textbf{Old Array (Capacity = 4):}};
    \node[cell] (o0) at (0, 1) {E1};
    \node[cell] (o1) at (0.8, 1) {E2};
    \node[cell] (o2) at (1.6, 1) {E3};
    \node[cell] (o3) at (2.4, 1) {E4};
    
    % New Array
    \node[anchor=west] (n_title) at (-4.8, -1.0) {\textbf{New Array (Capacity = 8):}};
    \node[moved] (n0) at (0, -1) {E1};
    \node[moved] (n1) at (0.8, -1) {E2};
    \node[moved] (n2) at (1.6, -1) {E3};
    \node[moved] (n3) at (2.4, -1) {E4};
    \node[cell] (n4) at (3.2, -1) {E5};
    \node[empty] (n5) at (4.0, -1) {};
    \node[empty] (n6) at (4.8, -1) {};
    \node[empty] (n7) at (5.6, -1) {};
    
    % Copying Arrows
    \draw[arrow, red, dashed] (o0) -- (n0);
    \draw[arrow, red, dashed] (o1) -- (n1);
    \draw[arrow, red, dashed] (o2) -- (n2);
    \draw[arrow, red, dashed] (o3) -- (n3);
    
    % Annotation
    \node[anchor=west, align=left, font=\small] at (6.4, 0.0) {
        1. Allocate capacity 8 \\
        2. Copy E1--E4 (red dashed) \\
        3. Insert new element E5 \\
        4. Free old memory \\
        5. Wasted cells: 3 (gray)
    };
\end{tikzpicture}
\caption{Visual Representation of the Dynamic Array Resizing Process}
\label{fig:array_resize_process}
\end{figure}

To ensure that the amortized time complexity of an insertion remains $O(1)$, 
conventional dynamic arrays utilize a geometric expansion policy. That is, the 
new capacity is determined by multiplying the current capacity by a constant 
growth factor, typically denoted as $g$. For example, the C++ Standard Template 
Library (STL) commonly employs a growth factor of $g = 2.0$ (doubling the 
capacity), while some implementations like Java's \texttt{ArrayList} use $g = 1.5$. 
Python's \texttt{list} implementation utilizes a size-dependent bit-shift 
formula:
\begin{equation}
    C_{\text{new}} = C_{\text{old}} + (C_{\text{old}} \gg 3) + (C_{\text{old}} < 9 ? 3 : 6)
\end{equation}
which roughly equates to an asymptotic growth factor of approximately $1.125$ 
for larger arrays.

Mathematically, the amortized cost can be proved using the \textit{Accounting Method}. 
Let the cost of a regular insertion without resize be $1$ unit of work (the 
actual write). Let the cost of copying an element during resize also be $1$ unit 
of work. If we use a growth factor of $2.0$, we can charge each insertion a 
``tax'' or amortized cost of $3$ units. When an element is inserted, $1$ unit pays 
for its immediate insertion. The remaining $2$ units are stored as credit on the 
element. When the array size reaches $N$ (where $N$ is a power of 2) and must be 
resized to $2N$, there are $N/2$ newly inserted elements since the last resize. 
Each of these $N/2$ elements has $2$ units of credit, totaling $N$ units of credit. 
This credit is exactly sufficient to pay for copying all $N$ elements (both the 
older $N/2$ elements and the newer $N/2$ elements) into the new array block. 
Thus, the amortized cost per insertion remains bounded by a constant, $O(1)$.

While geometric expansion guarantees amortized constant time for insertions over 
the long run, it introduces a rigid trade-off between memory waste and the 
frequency of reallocation. A high growth factor (e.g., $2.0$) minimizes the 
number of reallocations and the total number of elements copied, resulting in 
higher execution speed, but it can lead to substantial memory waste. For instance, 
if an array doubles its capacity to 100,000 elements but only 50,001 elements 
are ultimately inserted, nearly 50\% of the allocated memory remains unutilized. 
Conversely, a low growth factor (e.g., $1.125$) minimizes memory waste but forces 
the system to undergo frequent reallocations, drastically increasing the 
cumulative number of elements copied and incurring heavy performance penalties, 
especially for large arrays. 

This static, one-size-fits-all approach to resizing has remained largely 
unchanged for decades, despite the growing complexity and diversity of data 
ingestion patterns in modern applications.

\section{Problem Statement / Hypothesis}
Despite the widespread reliance on dynamic arrays, their conventional 
implementations utilize static, workload-independent resizing policies that 
fail to adapt to runtime data insertion patterns. Whether an application is 
slowly and predictably appending elements or suddenly experiencing a massive 
burst of incoming data, standard dynamic arrays react identically, constrained 
by their predefined growth factors (such as $1.5\times$ or $2.0\times$).

This lack of runtime adaptability gives rise to three main problem areas:
\begin{enumerate}
    \item \textbf{Suboptimal Memory Utilization:} Static policies often result 
    in severe memory waste, particularly at the end of an insertion phase. If 
    a large array expands right before the insertion process terminates, the 
    overarching memory waste can reach up to 50\% for a doubling policy. In large 
    data systems, this waste aggregates to gigabytes of unused memory.
    \item \textbf{Excessive Computational Overhead:} Conservative policies (e.g., 
    Python's $1.125\times$) that try to save memory often trigger excessive 
    reallocation events. The repetitive allocation of memory and the subsequent 
    copying of increasingly large numbers of elements severely degrade application 
    throughput and increase CPU utilization.
    \item \textbf{Lack of Workload Awareness:} Existing algorithms operate in a 
    vacuum, ignoring valuable runtime feedback such as the speed at which 
    elements are being inserted (insertion velocity) or the history of previous 
    reallocations. 
\end{enumerate}

\textbf{Research Hypothesis:} \\
\textit{By utilizing runtime feedback—specifically the resize history counter 
$k$ and the instantaneous insertion velocity $v$—a dynamic array growth algorithm 
can achieve a significantly better trade-off between memory waste and element-copying 
overhead than standard static policies. Specifically, the growth factor can be 
scaled down as the array matures (increasing $k$) to prevent massive memory waste, 
while scaling up during high-velocity bursts (increasing $v$) to minimize the 
frequency of costly reallocations.}

\section{Proposed Model}
To overcome the limitations of static growth factors, this study proposes a 
novel, workload-aware dynamic array growth algorithm. The fundamental innovation 
of the proposed model is the transition from a constant growth factor $g$ to a 
dynamic, continuously evaluated function $g(k, v)$, where the growth factor is 
determined at runtime based on two primary feedback variables:
\begin{enumerate}
    \item \textbf{Resize History ($k$):} The number of times the array has been 
    resized. This serves as a proxy for the current scale and maturity of the 
    array. As the array grows larger (higher $k$), the absolute cost of copying 
    elements increases, necessitating a carefully managed growth strategy to 
    prevent memory waste without triggering excessive copies.
    \item \textbf{Insertion Velocity ($v$):} The rate at which elements are being 
    appended to the array, calculated as $v = \frac{\Delta N}{\Delta t}$, where 
    $\Delta N$ is the number of elements inserted and $\Delta t$ is the elapsed 
    time since the last resize. This metric captures the real-time intensity of 
    the workload.
\end{enumerate}

To mathematically map these inputs to an optimal growth factor, this research 
systematically explores and develops three progressive architectural frameworks:
\begin{itemize}
    \item \textbf{Logarithmic Model (Formula A):} A model that dampens the growth 
    factor logarithmically as the array size increases, preventing aggressive 
    over-allocation at large capacities. The mathematical formulation is:
    \begin{equation}
        g(C) = 1 + \frac{\alpha}{\ln(C + \beta)}
    \end{equation}
    where $C$ is the current capacity, and $\alpha$ and $\beta$ are parameters 
    that control the rate of decay.
    
    \item \textbf{Harmonic Model (Formula B):} A model utilizing the harmonic 
    series, specifically designed to decay the growth factor as a function of 
    the resize count $k$:
    \begin{equation}
        g(k) = 1 + \frac{c_0}{k}
    \end{equation}
    where $k$ is the resize counter and $c_0$ is a scaling constant. This inherently 
    scales back expansion as the array matures.
    
    \item \textbf{Hybrid Framework (Formula C \& D):} The culmination of the 
    research, which synthesizes the strengths of both historical decay and 
    velocity-based scaling. The hybrid framework models the growth factor using 
    a workload-aware harmonic function: 
    \begin{equation}
        g(k,v) = \min\left(g_{\max}, \max\left(g_{\min}, 1 + \frac{c_0 + \lambda \ln(1+v)}{k}\right)\right)
    \end{equation}
    This final model is further optimized using Bayesian Optimization (GP-Optimized 
    Hybrid Framework) to fine-tune its hyper-parameters ($c_0, \lambda$) against 
    a unified objective score $J$ that penalizes both memory waste and computational 
    overhead (element copies).
\end{itemize}

\section{Objectives}

\subsection{Main Objective}
The principal objective of this study is to design, implement, and evaluate a 
workload-aware dynamic array growth algorithm that optimally balances memory 
efficiency and computational throughput by dynamically adapting its growth 
factor in response to runtime insertion patterns.

\subsection{Sub Objectives}
To achieve the main objective, the following sub-objectives have been established:
\begin{enumerate}
    \item To analyze and quantify the limitations, performance bottlenecks, and 
    memory inefficiencies of existing static resizing policies (e.g., $1.5\times$, 
    $2.0\times$, and Python's $1.125\times$) under diverse workload conditions.
    \item To mathematically formulate adaptive growth models, specifically 
    exploring Logarithmic and Harmonic functions, that incorporate resize history 
    and insertion velocity as feedback variables.
    \item To systematically optimize the hyper-parameters of the proposed models 
    using Bayesian Optimization (Gaussian Processes) to minimize a composite cost 
    function representing memory waste, element-copying overhead, and execution 
    time.
    \item To rigorously benchmark the proposed optimized Hybrid framework against 
    conventional static policies across multiple, highly varied workload profiles 
    (including uniform, bursty, and mixed patterns) spanning up to $10^7$ insertions.
\end{enumerate}

\section{Structure of the Thesis}
This dissertation is organized into five subsequent chapters to provide a 
comprehensive exploration of the research problem, the proposed solution, and the 
empirical findings. 
\begin{itemize}
    \item \textbf{Chapter 1: Introduction} outlines the context, identifies the 
    limitations of conventional dynamic arrays, formulates the problem statement, 
    and outlines the research objectives and proposed workload-aware models.
    \item \textbf{Chapter 2: Literature Review} provides an in-depth critical 
    analysis of existing literature surrounding memory allocation, dynamic data 
    structures, and algorithmic optimizations, establishing the foundation for 
    the proposed methodologies.
    \item \textbf{Chapter 3: Methodology} details the mathematical formulation of 
    the adaptive growth functions, the design of the workload profiles, the 
    architecture of the custom benchmarking suite implemented in Python, and the 
    Bayesian Optimization pipeline utilized for parameter tuning.
    \item \textbf{Chapter 4: Results and Discussion} presents the empirical 
    evaluation of the proposed models against static baselines, analyzing metrics 
    such as memory waste, reallocation frequency, cumulative element copies, and 
    execution throughput across various scales and workload patterns.
    \item \textbf{Chapter 5: Conclusion} summarizes the key findings of the 
    research, highlights the contributions to the field of software engineering, 
    discusses the limitations of the current study, and proposes directions for 
    future work and real-world deployment.
\end{itemize}

\section{Summary for the Chapter}
This chapter established the foundational context for the research by highlighting 
the ubiquitous nature of dynamic arrays and the critical trade-off between memory 
efficiency and computational throughput governed by their growth strategies. The 
inherent limitations of conventional, static resizing policies were identified as 
the primary catalyst for this study, particularly their inability to adapt to 
varying data insertion workloads. To address these inefficiencies, an overview of 
a novel, workload-aware adaptive growth algorithm was presented, combining resize 
history and real-time insertion velocity. Finally, the main and specific objectives 
were clearly delineated, laying the groundwork for the mathematical formulation, 
implementation, and rigorous empirical evaluation detailed in the subsequent 
chapters of this thesis.
